\documentclass{article}
\usepackage{tcee-repaso}
\tcetema[Finanzas y presupuesto de la UE]{3.B.37}{Las finanzas de la Unión Europea y el presupuesto comunitario. El marco financiero plurianual actual.}
\tcefecha{19/11/2026}
\begin{document}
\cabecerarepaso
\begin{columnas}

% ==========================================================================
\bloque[2,5]{Introducción}
\begin{itemize}
\item \textbf{Enganche}: «Cuentas claras, amistades largas».
  \begin{itemize}
  \item La Comisión publica cuánto paga y recibe cada Estado.
  \item Aun así, el MFP 2028-2034 (casi 2~billones, 1,26\,\% de la RNB) sigue sin acuerdo en noviembre de 2026.
  \end{itemize}
\item \textbf{Relevancia}: instrumento de las políticas comunes; gran emisor de deuda ($\approx$875.000~M€); España, 77.600~M€ del MRR \vertema{4.A.30}.
\item \textbf{Contextualización}: \concepto{presupuesto} (art.~310~TFUE) dentro del \concepto{MFP} (7~años).
  \begin{itemize}
  \item Sin poder tributario ni déficit; $\approx$1\,\% de la RNB (Estados $\approx$50\,\%; EE.~UU. federal $\approx$25\,\%).
  \item Concentrado: RNB $\approx$2/3 de los ingresos; PAC + cohesión $\approx$2/3 del gasto.
  \end{itemize}
\item \textbf{Problemática}: ¿qué reglas? ¿Cómo se financia? ¿En qué gasta y debería gastar? ¿Puede endeudarse?
\item \textbf{Estructura}: \hyperref[rep:I]{I} reglas; \hyperref[rep:II]{II} cifras; \hyperref[rep:III]{III} fuera del presupuesto.
\end{itemize}
\esquemacante{%
\apartadocante{Introducción}{2,5}
\apartadocante{I. La constitución financiera}{8,5}
\subcante{1. Los textos}{1}
\subcante{2. Principios y ciclo presupuestario}{2,5}
\subcante{3. Ingresos: recursos propios}{2,5}
\subcante{4. Gastos: el MFP}{2,5}
\apartadocante{II. Aproximación descriptiva}{9}
\subcante{1. Evolución histórica}{2,5}
\subcante{2. MFP 2021-2027, 2026 y saldos}{4,5}
\subcante{3. La propuesta 2028-2034}{2}
\apartadocante{III. Actividad extrapresupuestaria}{9}
\subcante{1. Instrumentos especiales}{1}
\subcante{2. Préstamos y endeudamiento}{3}
\subcante{3. NextGenerationEU}{5}
\apartadocante{Conclusión}{1}}

% ==========================================================================
\bloque[8,5]{I. La constitución financiera\label{rep:I}}
\begin{ideaclave}
Sin poder tributario ni déficit: los Estados fijan por unanimidad cuánto ingresa (DRP) y cuánto gasta en 7~años (MFP); dentro, el PE y el Consejo aprueban cada año el presupuesto.
\end{ideaclave}

\sub{I.1. Los textos de la constitución financiera\label{rep:I.1}}
\begin{itemize}
\item \emph{Tratados}: arts.~310-325 TFUE; 41~TUE.
\item \emph{Reglamento financiero} 2024/2509: PLO (art.~322); aplicable desde 30-9-2024.
\item \emph{Decisión de recursos propios} 2020/2053: unanimidad, consulta al PE, ratificación de los 27 (art.~311).
\item \emph{Reglamento del MFP} 2020/2093 (revisado por 2024/765): unanimidad + aprobación del PE (art.~312).
  \begin{itemize}
  \item Acuerdo interinstitucional (16-12-2020); condicionalidad del Estado de Derecho (2020/2092).
  \end{itemize}
\end{itemize}
\graficorepaso[0.78\linewidth]{constitucion-financiera}{Cada texto encaja en el de arriba}

\sub{I.2. Principios y ciclo presupuestario\label{rep:I.2}}
\subsub{a) Los principios presupuestarios}
\begin{itemize}
\item 9~principios:
  \begin{itemize}
  \item unidad y veracidad; equilibrio; unidad de cuenta;
  \item anualidad: créditos de compromiso y de pago (RAL 501.100~M€ en 2025);
  \item universalidad: no afectación ni compensación (excepción: ingresos afectados externos, NGEU);
  \item especialidad; transparencia; buena gestión financiera.
  \end{itemize}
\end{itemize}
\subsub{b) El ciclo presupuestario}
\subsubsub{Elaboración y aprobación (art.~314)}
\begin{itemize}
\item Proyecto (1-IX; junio en la práctica) $\to$ Consejo (1-X) $\to$ PE (42~días) $\to$ Consejo (10~días) $\to$ Conciliación (21) $\to$ texto conjunto (14).
\item El PE tiene la última palabra. Sin presupuesto: doceavas partes (art.~315) \vertema{3.B.36}.
\end{itemize}
\graficorepaso[0.9\linewidth]{procedimiento-presupuestario}{Art.~314: una lectura en cada institución y conciliación}
\subsubsub{Ejecución (art.~317)}
\begin{itemize}
\item Compartida $\approx$70\,\%, directa $\approx$20\,\%, indirecta $\approx$10\,\%.
\end{itemize}
\subsubsub{Control}
\begin{itemize}
\item Tribunal de Cuentas (2025): opinión desfavorable sobre el gasto, error 3,8\,\%; MRR con reservas.
\item Descargo del PE ($t-2$, art.~319); OLAF.
\item Fiscalía Europea: 3.602~investigaciones, 67.270~M€ de perjuicio.
\item Condicionalidad: Hungría, 6.300~M€ (2022); levantada en 2026.
\end{itemize}

\sub{I.3. Ingresos: la Decisión de recursos propios (2020)\label{rep:I.3}}
\subsub{a) El sistema de recursos propios}
\begin{itemize}
\item \emph{Techo}: \concepto{1,40\,\%} de la RNB (pagos), 1,46\,\% (compromisos) + 0,6~temporal para el NGEU (hasta 2058).
  \begin{itemize}
  \item El margen garantiza la deuda.
  \end{itemize}
\item \emph{Recursos} (2025, \% de los recursos propios):
  \begin{itemize}
  \item tradicionales (aduanas; retención del 25\,\%): 15\,\%;
  \item IVA (0,30\,\% de una base $\leq$50\,\% de la RNB): 15\,\%;
  \item RNB (de cierre, desde 1988): 65\,\% (70\,\% en 2026);
  \item plástico (0,80 €/kg, desde 2021): 4\,\%.
  \end{itemize}
\item \emph{Correcciones}: cheque británico (1984, 2/3); 2021-2027, a tanto alzado: DK, DE, NL, AT, SE.
\end{itemize}
\graficorepaso[0.9\linewidth]{ingresos-composicion}{De contribuciones y aranceles al IVA y, desde los 90, a la RNB}
\subsub{b) Los nuevos recursos propios}
\begin{itemize}
\item Propuestas: 2021 (RCDE 25\,\%, MAFC 75\,\%, Pilar 1); 2023 (RCDE 30\,\%, 0,5\,\% de los beneficios); 2025 (5~recursos, 58.500~M€/año).
  \begin{itemize}
  \item Ninguno aprobado (nov. 2026).
  \end{itemize}
\item \concepto{MAFC} \vertema{3.B.31}: el importador de cemento, acero, aluminio, fertilizantes, electricidad e hidrógeno compra certificados al precio del RCDE.
  \begin{itemize}
  \item Por el CO\textsubscript{2} incorporado, menos el precio pagado en origen y los derechos gratuitos (97,5\,\% en 2026, 0~en 2034).
  \item Definitivo desde 1-1-2026; certificados desde II-2027; exentos $<$50~t/año.
  \end{itemize}
\end{itemize}
\graficorepaso[0.78\linewidth]{mafc}{El MAFC: el importador paga el mismo precio del carbono}

\sub{I.4. Gastos: el marco financiero plurianual\label{rep:I.4}}
\subsub{a) Qué es y cómo se aprueba}
\begin{itemize}
\item \concepto{MFP}: techos anuales de compromiso por rúbricas y de pago, $\geq$5~años.
  \begin{itemize}
  \item Previsible, disciplina; concentra el conflicto.
  \item Sin acuerdo, se prorroga el último año (art.~312.4).
  \end{itemize}
\end{itemize}
\subsub{b) ¿Qué debe financiar la Unión? Los bienes públicos europeos}
\begin{itemize}
\item \emph{Idea}: federalismo fiscal (\autor{Musgrave}, 1959; \autor{Oates}, 1972).
  \begin{itemize}
  \item Si los beneficios cruzan fronteras, cada Estado ofrece menos de lo eficiente (polizón).
  \end{itemize}
\item \emph{Resultado}: regla de \autor{Bowen} (1943)-\autor{Samuelson} (1954): suma de lo que valoran todos = coste.
  \begin{itemize}
  \item La cumple el presupuesto común; la brecha crece con el n.º de Estados: \emph{valor añadido europeo}.
  \end{itemize}
\item \emph{Implicaciones}: I+D, redes, fronteras, defensa, clima.
\item \emph{Evidencia}: \autor{Alesina}-\autor{Angeloni}-\autor{Schuknecht} (2005); \autor{Draghi} (2024): 750.000-800.000 M€/año.
\item \emph{Valoración}: supone un planificador; deciden 27~por unanimidad.
\end{itemize}

% ==========================================================================
\bloque[9]{II. Aproximación descriptiva\label{rep:II}}
\begin{ideaclave}
Pequeño ($\approx$1\,\% de la RNB) y cambiado por dentro: de aranceles a RNB; de la PAC a cohesión y competitividad. Cada cambio, por unanimidad y con compensaciones.
\end{ideaclave}

\sub{II.1. Evolución histórica\label{rep:II.1}}
\subsub{a) De las contribuciones a los recursos propios (1951-1987)}
\begin{itemize}
\item Contribuciones (la CECA, con exacción propia) $\to$ PAC (FEOGA 1962; $>$90\,\% agrario en 1970).
\item Recursos propios (Decisión de 21-4-1970; IVA cobrado desde 1979).
\item PE con poder presupuestario (Tratados de 1970 y 1975); Tribunal de Cuentas (1977).
\item Crisis de los 80: presupuestos rechazados (1980, 1985); Fontainebleau (1984: cheque, IVA 1,4\,\%).
\end{itemize}
\subsub{b) La reforma de 1988 y los primeros marcos (1988-2006)}
\begin{itemize}
\item \emph{Delors I}: RNB, techo 1,15-1,20\,\%, perspectivas financieras, PAC $\leq$74\,\% del crecimiento del PNB, fondos estructurales $\times2$.
\item \emph{Delors II} (Edimburgo 1992): techo 1,27\,\%; IVA al 1\,\% y base $\leq$50\,\%; Fondo de Cohesión (1993).
\item \emph{Agenda 2000} (Berlín 1999): IVA 0,50\,\%; retención 25\,\%; preadhesión.
\end{itemize}
\subsub{c) La Unión ampliada y en crisis (2007-2027)}
\begin{itemize}
\item 2007-2013: el Este; IVA 0,30\,\%; cohesión 35\,\%.
\item 2014-2020: 1\,\% de la RNB, primer recorte real; Lisboa: el MFP, reglamento.
\item 2021-2027: NGEU, plástico, 1,40\,\%, fin del cheque.
\item Gasto $\approx$1\,\% de la RNB desde 1988 (máx. 1,13\,\% en 1993; 1,57\,\% en 2021 con NGEU).
\end{itemize}
\graficorepaso[0.78\linewidth]{gastos-evolucion}{Pagos por áreas desde 1958 y \% de la RNB: casi plano desde 1988 hasta el NGEU}

\sub{II.2. El MFP 2021-2027 y el presupuesto de 2026\label{rep:II.2}}
\subsub{a) El MFP 2021-2027 y NextGenerationEU}
\begin{itemize}
\item Consejo Europeo 17-21 de julio de 2020; Reglamento 17-12-2020.
\item 1.074.300~M€ de 2018 (1.225.900~corrientes tras 2024; 1,02\,\% RNB) + NGEU 750.000: $\approx$2~billones.
\item \emph{Rúbricas} (M€ de 2018):
  \begin{itemize}
  \item mercado único 152.000 (12\,\%); cohesión 432.200 (35\,\%) + NGEU 776.500;
  \item recursos naturales 400.300 (33\,\%); migración 28.900; seguridad y defensa 16.400;
  \item vecindad 113.700 (9\,\%); administración 82.500 (7\,\%).
  \end{itemize}
\item Clima $\geq$30\,\%.
\end{itemize}
\graficorepaso[0.78\linewidth]{mfp-2021-2027}{MFP y NGEU por rúbricas (miles de millones de euros)}
\subsub{b) La revisión intermedia de 2024}
\begin{itemize}
\item Reglamento 2024/765: +64.600 M€.
  \begin{itemize}
  \item Mecanismo para Ucrania 50.000 (17.000~subvenciones + 33.000~préstamos);
  \item STEP; migración; EURI (intereses del NGEU).
  \end{itemize}
\end{itemize}
\subsub{c) El presupuesto de 2026}
\begin{itemize}
\item 192.768~M€ en compromisos, 190.092~en pagos.
\item Ingresos: RNB 70\,\%, IVA 13\,\%, tradicionales 11\,\%, plástico 4\,\%.
\item 2027: acuerdo en conciliación (10-11-2026): 199.708 / 211.520~M€.
\end{itemize}
\subsub{d) Contribuyentes y receptores netos}
\begin{itemize}
\item Saldos 2025 (sin NGEU): Alemania $-$14.855~M€ ($-$0,32\,\% RNB); Polonia +6.674; Letonia +2,7\,\%.
\item \emph{España}: $-$0,06\,\% (+0,38\,\% con NGEU).
\end{itemize}
\graficorepaso[0.78\linewidth]{saldos-netos}{Saldos operativos 2025 (\% RNB): barras sin NGEU, círculos con él}
\subsub{e) La lógica del justo retorno}
\begin{itemize}
\item \emph{Idea} (\autor{Monti}, 2016): cada gobierno mira su saldo neto.
  \begin{itemize}
  \item El gasto asignable es suma cero y se ve.
  \item El bien público beneficia a todos de forma difusa; solo se ve lo que cuesta $\Rightarrow$ nadie lo defiende.
  \end{itemize}
\item \emph{Consecuencias}: sesgo al gasto asignable, correcciones, sobres, inercia (prórroga).
\item \emph{Evidencia}: \autor{Kauppi}-\autor{Widgrén} (2004): el poder de voto explica el reparto.
\item \emph{Salidas}: bienes públicos, recursos propios genuinos, deuda para crisis.
\end{itemize}

\sub{II.3. La propuesta 2028-2034\label{rep:II.3}}
\begin{itemize}
\item 16-7-2025: 1,763~billones de 2025 ($\approx$1,98~corrientes; 1,26\,\% RNB, 0,11~para el NGEU).
  \begin{itemize}
  \item Planes nacionales y regionales 865.000 (PAC + cohesión);
  \item Fondo de Competitividad 409.000 (+ Horizonte 175.000);
  \item Europa Global 200.000; Reserva para Ucrania; crisis: préstamos hasta 395.000.
  \end{itemize}
\item 5~nuevos recursos (RCDE, MAFC, residuos electrónicos, tabaco, CORE): 58.500~M€/año; techo 1,75\,\%.
\end{itemize}
\graficorepaso[0.78\linewidth]{mfp-2028-2034}{La propuesta 2028-2034: 4~bloques y 5~nuevos recursos}
\begin{itemize}
\item \emph{Negociación}:
  \begin{itemize}
  \item PE (abril de 2026): +10\,\%; caja irlandesa (octubre): 1,622~billones ($-$8\,\%);
  \item frugales, más recortes; España, con los Amigos de la Cohesión;
  \item cumbre extraordinaria, 26-27 de noviembre de 2026.
  \end{itemize}
\end{itemize}

% ==========================================================================
\bloque[9]{III. Actividad extrapresupuestaria\label{rep:III}}
\begin{ideaclave}
No puede endeudarse para el presupuesto, pero sí para prestar y, con autorización de la DRP, para gastar: garantía, el margen del techo. Deuda de la UE: $\approx$875.000~M€.
\end{ideaclave}

\sub{III.1. Instrumentos especiales y fondos fuera\label{rep:III.1}}
\begin{itemize}
\item \emph{Por encima de los techos}: Flexibilidad, RSAE (Fondo de Solidaridad), FEAG, Reserva del Brexit, Reserva para Ucrania, EURI.
\item \emph{Fuera}: FEAP (art.~41.2~TUE), fondos fiduciarios; el FED, integrado en 2021.
\end{itemize}

\sub{III.2. Préstamos, endeudamiento y estabilización\label{rep:III.2}}
\subsub{a) ¿Puede endeudarse la Unión?}
\begin{itemize}
\item Bases: arts.~122, 143, 212, 352; préstamos espejo; financiación unificada desde 2023.
\end{itemize}
\subsub{b) Préstamos sectoriales y macroeconómicos}
\begin{itemize}
\item CECA y Euratom; balanza de pagos (art.~143, 50.000~M€); MEEF (2010, 60.000).
\item FEEF y MEDE (intergubernamentales; 500.000; España 41.300) \vertema{3.B.45}.
\item SURE (98.400; España 21.324); SAFE (150.000, defensa; España 1.000).
\item Ucrania: AMF+ 18.000, G7 18.100, Mecanismo 33.000, Préstamo de Apoyo 90.000 (2026-2027).
\end{itemize}
\subsub{c) El Banco Europeo de Inversiones}
\begin{itemize}
\item InvestEU: 26.200~M€ de garantía \vertema{3.B.30}.
\end{itemize}

\sub{III.3. El Fondo NextGenerationEU\label{rep:III.3}}
\subsub{a) Idea y diseño}
\begin{itemize}
\item 6~rasgos:
  \begin{itemize}
  \item deuda de la UE (750.000~de 2018), devolución 2028-2058 (máx. 7,5\,\% de 390.000~al año);
  \item garantía: 0,6~de RNB del techo;
  \item subvenciones 390.000 + préstamos 360.000;
  \item gastan los Estados (MRR, 90\,\%); condicionalidad (hitos y objetivos).
  \end{itemize}
\item Base: art.~122; ingresos afectados externos.
\end{itemize}
\subsub{b) El Mecanismo de Recuperación y Resiliencia}
\begin{itemize}
\item Reparto 70/30; planes con $\geq$37\,\% clima y $\geq$20\,\% digital; evaluación A/B/C; freno de emergencia.
\item Final: hitos, 31-8-2026; pagos, 31-12-2026.
\item 1-10-2026: 450.000~M€ pagados (79\,\%), 68\,\% de hitos.
\end{itemize}
\graficorepaso[0.78\linewidth]{mrr-paises}{MRR recibido 2020-2025 (\% del PIB): España, 4,2\,\%}
\subsub{c) La financiación: la Comisión en los mercados}
\begin{itemize}
\item 836.800~M€ de deuda (6-2026); 180.000~en 2026; coste 3,32\,\%; AAA (S\&P AA+).
\end{itemize}
\subsub{d) España y el MRR}
\begin{itemize}
\item PRTR + adenda (2023): 163.000~M€; recorte de préstamos (dic. 2025): plan de 102.000.
\item 6~pagos, 77.600~M€; 7.º (25.861) pedido el 30-9-2026.
\end{itemize}
\subsub{e) ¿Es sostenible la deuda común?}
\begin{itemize}
\item \emph{Modelo} (\autor{Domar}, 1944) \vertema{3.A.39}: importa la deuda sobre la renta; con $i<\gamma$ baja sola.
  \begin{itemize}
  \item Calendario cerrado (solo la deuda de las subvenciones, 31~años), pagado con recursos propios.
  \item $D_0\approx421{.}000$ M€, $i=3{,}3$\,\%, $\gamma=3{,}5$\,\%: 27.500~M€ (0,14\,\% RNB) en 2028 $\to$ 0,03\,\% en 2058.
  \end{itemize}
\item \emph{Valoración}: sostenible; la duda es quién paga (nuevos recursos, recortes o refinanciar).
  \begin{itemize}
  \item \autor{Hamilton} (1790), \autor{Sargent} (2012): deuda común con impuestos comunes; activo seguro.
  \item \autor{Becker} et al. (2010): +1,6 pp de crecimiento en las regiones de objetivo 1.
  \end{itemize}
\end{itemize}
\graficorepaso[0.78\linewidth]{deuda-ngeu}{La deuda baja sola; la carga es pequeña y decreciente}

% ==========================================================================
\bloque[1]{Conclusión}
\begin{itemize}
\item \textbf{Recapitulación}:
  \begin{itemize}
  \item constitución financiera rígida;
  \item presupuesto pequeño, financiado por RNB, 2/3 en PAC y cohesión (unanimidad y justo retorno);
  \item fuera, la UE, gran emisor de deuda.
  \end{itemize}
\item \textbf{Relevancia}: España, mayor transferencia de su historia (MRR); saldo ordinario cerca del equilibrio.
\item \textbf{Extensiones}: \vertema{3.B.38} \vertema{3.B.40} \vertema{3.B.42} \vertema{3.B.45} \vertema{4.B.14}.
\item \textbf{Opinión}: hace bien lo que hace, pero poco de lo que debería (bienes públicos europeos).
  \begin{itemize}
  \item La propuesta 2028-2034 va bien si no acaba en un reparto de sobres; método \autor{Schuman}: «realizaciones concretas» y «solidaridad de hecho».
  \end{itemize}
\item \textbf{Idea final}: las cuentas miden lo que vuelve a cada socio; la amistad será larga si se mira lo que solo se consigue juntos.
\end{itemize}

% ==========================================================================
\anexorepaso
\bloque{Datos clave}
\begin{itemize}
\item MFP 2021-2027: \dato{1.074.300 M€ de 2018}{1,02\,\% RNB} + NGEU 750.000.
\item Presupuesto 2026: \dato{192.768 M€}{compromisos}.
\item Ingresos 2026: \dato{RNB 70\,\%, IVA 13\,\%, RPT 11\,\%}{Comisión}.
\item MRR: \dato{450.000 M€ pagados (79\,\%)}{1-10-2026}; España \dato{77.600 M€}{6 pagos}.
\item Deuda de la UE: \dato{$\approx$875.000~M€}{2026}.
\end{itemize}
\bloque{Autores y fechas}
\hito{1944}{\autor{Domar}: sostenibilidad de la deuda}
\hito{1954}{\autor{Samuelson}: bienes públicos}
\hito{1970}{Decisión de recursos propios}
\hito{1972}{\autor{Oates}: federalismo fiscal}
\hito{1984}{Fontainebleau: cheque británico}
\hito{1988}{Delors I: recurso RNB}
\hito{2020}{MFP 2021-2027 y NGEU}
\hito{2025}{Propuesta 2028-2034}
\bloque{Preguntas del tribunal}
\begin{enumerate}
\item \textbf{¿Puede endeudarse la UE?} No para el presupuesto (art.~310); sí para prestar y, con la DRP, para gastar.
\item \textbf{¿Créditos de compromiso y de pago?} Obligaciones del año frente a pagos del año; RAL.
\item \textbf{¿Por qué el recurso RNB es el principal?} De cierre, proporcional a la renta; el IVA es regresivo y los aranceles caen.
\item \textbf{¿Qué pasa si no hay MFP o presupuesto?} Prórroga del último año (312.4); doceavas partes (315).
\item \textbf{¿NGEU y art.~125?} No asume deudas de los Estados: deuda propia para un proyecto común.
\item \textbf{¿Qué es el margen de maniobra?} Techo de recursos propios menos gasto: garantía de la deuda.
\end{enumerate}
\begin{trampa}
\begin{itemize}
\item Techo: \emph{1,40\,\%} (pagos), \emph{1,46\,\%} (compromisos) + \emph{0,6}.
\item MFP: unanimidad + \emph{aprobación} del PE; DRP: \emph{consulta} y ratificación.
\item Descargo: ejercicio \emph{$t-2$}; informe del TCE: \emph{$t-1$}.
\item IVA: \emph{0,30\,\%}, base $\leq$\emph{50\,\%}; plástico: \emph{0,80} €/kg.
\item NGEU: \emph{750.000} M€ de \emph{2018}; devolución \emph{2028-2058}.
\end{itemize}
\end{trampa}

\end{columnas}
\end{document}
